\subsection{经期性行为的常见疑问速答}

本章前面已给出"可以，但需双方自愿与注意卫生"的基本答案。这里把实际被问得最多的几个问题集中回答一次，便于快速查阅。

\begin{itemize}
\item \textbf{经期真的不会怀孕吗？}——\textbf{会，不是安全期}。精子在女性生殖道内可存活 3--5 天，而排卵时刻受周期波动影响；周期较短（如 21--24 天）或经期较长者，经期结束时可能已接近排卵，无保护性交仍有受孕可能；
\item \textbf{经期做会不会得病？}——不是必然，但风险确实略升高。宫颈口在经期微开，经血是细菌的良好培养基，女性的盆腔感染风险略增；对男性而言，经血中的病原体（如 HIV、乙肝病毒）可经尿道口黏膜形成暴露。\textbf{安全套能同时降低这两个方向的风险}，这也是经期使用安全套最主要的价值；
\item \textbf{一定很麻烦、很脏吗？}——有具体做法可以把它变得不麻烦：铺深色毛巾、选择浴室环境（清理方便）、优先选经量较少的时段。\textbf{"脏"是观念判断，不是卫生事实}；
\item \textbf{会不会加重子宫内膜异位症？}——目前\textbf{没有证据}表明经期性交会导致或加重内异症。经血逆流是内异症的主流发病学说之一，但逆流在几乎所有女性中都存在，而内异症只发生在少数人身上——理论成立不等于"经期性交导致内异症"这个推论成立（详见女性卷相关小节的完整讨论）；
\item \textbf{想用它缓解痛经，可以吗？}——可能有用：高潮时释放的内啡肽与催产素、子宫收缩后的放松期，都可能带来短时缓解，个体差异较大。但它是"可能有效"而非治疗方法，替代不了就医。若痛经每月严重影响学习与工作、止痛药无效，应排查子宫内膜异位症、腺肌症等继发原因；
\item \textbf{一方不愿意怎么办？}——\textbf{拒绝不需要理由}。不偏好经期性交的原因可能是观念、洁癖、身体不适或单纯没有兴致，这些都正常；反过来，"医学上可以做"也不构成说服对方同意的理由；
\item \textbf{哪些情况必须避免？}——流产后或产后的医嘱恢复期内、放置宫内节育器后的初期、盆腔炎急性期、有明确医嘱禁止的情形；若性交后出现下腹痛、发热，或经量骤增、行经期明显延长，应及时就诊；
\item \textbf{避孕措施能省吗？}——不能。经期既不是安全期，也不能预防性传播感染。
\end{itemize}

\begin{tcolorbox}[colback=pink!5!white,colframe=red!50!black,title=贴心小叮咛]
回答"经期能不能做"，需要同时说清两句话：\textbf{医学上没有绝对禁忌}，以及\textbf{你完全有权不做}。前一句防止人被羞耻感绑架，后一句防止人被"科学证明可以做"绑架。
\end{tcolorbox}

